\section{Main result}

% \bigskip

In this section, we characterize the values of $n$ for which the puzzle has a strategy and we present a sketch of the proof. The details of the proof are given in the next sections and are adapted for formalization in Lean.

\smallskip

\noindent {\bf Theorem}
{\it The puzzle has a strategy for $n\in\mathbb{N}^*$ if and only if $n$ is a power of $2$.}

\medskip

{\it Sketch of the proof.}

$\boxed{\text{Modus Ponens}}$ For every $b\in\Zz^n$,
$$\Fin n\ni i\mapsto f(b+{\rm e}_i)\in \Fin n \text{ is surjective},$$ and thus also injective.
Fix some $k \in \Fin n$ (e.g., $k=0 < n$).
Then
\begin{align*}
2^n & = \sum_{b\in\Zz^n} 1  = \sum_{b\in\Zz^n} \sum_{i\in\Fin n}
\delta(f(b+{\rm e}_i),k) = \sum_{i\in\Fin n} \sum_{b\in\Zz^n} \delta(f(b+{\rm e}_i),k)\\
& = \sum_{i\in\Fin n} \sum_{b\in\Zz^n} \delta(f(b),k)  = n \boldsymbol{\cdot} \sum_{b\in\Zz^n} \delta(f(b),k).
\end{align*}

% \bigskip

$\boxed{\text{Modus Ponens Reverse}}$  Let $n=2^m$. Let $g : \Fin (2^m) \to \Zz^m$ be a bijection. Define
$$f(b) := g^{-1}\left(\sum_{j\in\Fin 2^m} b_j \boldsymbol{\cdot} g(j)\right), \quad b\in\Zz^n.$$

Let $b\in\Zz^n$ and $k\in\Fin n$.
\begin{align*}
& f(b+{\rm e}_i) = k
 \Leftrightarrow \sum_{j\in\Fin 2^m} b_j \boldsymbol{\cdot} g(j) +
                  \sum_{j\in\Fin 2^m} {\rm e}_{ij} \boldsymbol{\cdot} g(j) = g(k) \\
& \Leftrightarrow \sum_{j\in\Fin 2^m} b_j \boldsymbol{\cdot} g(j) + g(i) = g(k)  \Leftrightarrow i = g^{-1}\left(g(k) - \sum_{j\in\Fin 2^m} b_j \boldsymbol{\cdot} g(j)\right).
\end{align*}
